% Preamble for bible.tex: page layout, fonts, and the project's macros.
% bible.tex itself holds only the book list.

% Page layout options
% Note: the "showframe" option is useful for debugging.
\usepackage[paperwidth=5.25in, paperheight=8in, top=0.5in, bottom=0.25in, outer=1in, inner=0.44in, marginparwidth=0.5in, headsep=0.125in, footskip=0.25in]{geometry}
% \setlength{\topmargin}{-0.7in} % Relative to 1" default % This is for line spacing or something. See Bringhurst.

% Font spacing options
% \renewcommand\baselinestretch{12pt} % Kills compilation. % For line spacing.
\textwidth=23pc
\parindent=2em

% Packages that don't necessarily belong elsewhere.
\usepackage{mathtools}
\usepackage{nicefrac}
\usepackage{textcomp} % For miscellaneous symbols, viz.: copyright, degree.

% Hyphen penalty (higher numbers discourage hyphenation, up to 10000).
% 4000 seems to work fine because it only allows hyphenation in extreme conditions (and doesn't require having to use ~ all the time). 8000 was previously being used but anything with a badness over 5000 is ugly to the naked eye, so 4000 is now being used (5000 wasn't enough because it allowed things in the range 4800--5200).
% Never use 10000.
\sloppy\hyphenpenalty=4000

% Widow and orphan penalty.
\widowpenalty=10000
\clubpenalty=10000

% %%%%%%%%%%%%%%%%%%
%     APPENDICES
% %%%%%%%%%%%%%%%%%%
\usepackage[toc,page]{appendix}

% %%%%%%%%%%%%%%%%%%%%%%%%%%%
%     CHAPTER and SECTION
% %%%%%%%%%%%%%%%%%%%%%%%%%%%
% The "explicit" option requires the title to be given, supplied as #1.
\usepackage[explicit]{titlesec}

% Note: options for reconfiguring titles are found in the document.

% Defining two new commands which will be dynamically changed on each call of \book
\newcommand\englishchaptertitle{}
\newcommand\hebrewchaptertitle{}

%% Shorter way of declaring a chapter with both a Hebrew and an English name.
\newcommand\book[2]{%
  \chapter[#1\hfill#2~~~]{#1}% ~ added for spacing in ToC.
  \pdfbookmark[0]{#1}{book.#1}% Top-level PDF bookmark for the book.
  \renewcommand\englishchaptertitle{#1}%
  \renewcommand\hebrewchaptertitle{#2}%
  \thispagestyle{fancy}% Forces no page number in cfoot.
  \markboth{1\thinspace:\thinspace1}{1\thinspace:\thinspace1}%
}

% %%%%%%%%%%%%%%%%%%%%%%%%
%     CHAPTER HEADINGS
% %%%%%%%%%%%%%%%%%%%%%%%%
\newcommand\heading[2]{%
  \medskip%
  \pdfbookmark[1]{\englishchaptertitle\ #1}{\englishchaptertitle.#1}% Chapter bookmark under its book.
  \markboth{#1:1}{#1:1}%
  \marginpar{%
    \vspace{-2em}%
    \color{Gray}%
    {%
      \Huge%
      \textbf{%
        \hspace*{0.1em}%
        \begin{center}%
          #1%
        \end{center}%
      }%
    }%
  }%
  \noindent%
  \small%
  \emph{%
    \newline%
    #2.%
  }%
  \smallskip%
}

\usepackage{epigraph}

% %%%%%%%%%%%%%%%%%%%%%%%%%%%
%     ENUMERATE and VERSE
% %%%%%%%%%%%%%%%%%%%%%%%%%%%
% paralist allows for separate paragraphs in the same enumeration environment.
\usepackage{paralist}

% Change verse number to red superscript.
\usepackage[usenames,dvipsnames,svgnames]{xcolor}
\renewcommand\labelenumi{\hspace*{0.5em}\color{Red}{$^{\arabic{enumi}}\hspace*{-0.1em}$}}

% Patch enumi counter to always start with 1.
\usepackage{etoolbox}
\newcounter{myc}[enumi]
\patchcmd{\enumi}{}{\setcounter{myc}{1}}{}{}

\newcommand\vrs[1]{%
  \markboth{#1}{#1}% Sets up header. Stores #1 as both \leftmark and \rightmark.
  \item% Print verse number.
}

% Allows for more control of enumerations and lists.
\usepackage{enumitem}

% %%%%%%%%%%%%%%%%%%%%%%%%%%%%%%%
%     FONTS and ABBREVIATIONS
% %%%%%%%%%%%%%%%%%%%%%%%%%%%%%%%
\usepackage{fontspec}
\usepackage{metalogo}

% Egyptian hieroglyphs
\usepackage{graphicx}

% Greek
% Project fonts load straight from fonts/, so nothing needs installing system-wide.
\newfontfamily\sblg[Path=fonts/,Script=Greek,Scale=MatchUppercase]{SBL_grk.ttf}
\newcommand\Greek[1]{{\sblg #1}}

% Hebrew
\newfontfamily\sblh[Path=fonts/,Script=Hebrew,Scale=MatchUppercase]{SBL_Hbrw.ttf}
\newcommand\Hebrew[1]{{\sblh #1}}

% Phoenician
\usepackage{phoenician}

% Syriac
% Noto Sans Syriac's default style is Estrangela.
\newfontfamily\sbls[Path=fonts/,Script=Syriac,Scale=MatchUppercase]{NotoSansSyriac-Regular.ttf}
\newcommand\Syriac[1]{{\sbls{}#1}}

% ABBREVIATIONS
\newcommand\lord{\textsc{Lord}}
\newcommand\god{\textsc{God}}

\newcommand\caref[3]{#1~#2,#3}
\newcommand\haref[3]{#1~#2$_\text{#3}$}
\newcommand\vref[3]{#1~#2\thinspace:\thinspace#3}

\newfontfamily\abbrev[Path=fonts/,Scale=MatchUppercase,BoldFont=AppSILB.TTF,ItalicFont=AppSILI.TTF,BoldItalicFont=AppSILBI.TTF]{AppSILR.ttf}
\newcommand\aquila{\Greek{α´}}
\newcommand\arabica{{\abbrev{}𝔄}}
\newcommand\bomberg{{\abbrev{}𝔅}}
\newcommand\hadagger{{$\dagger$}}
\newcommand\fragheb{{\abbrev{}ℭ}}
\newcommand\latina{{\abbrev{}𝔏}}
\newcommand\masoretic{{\abbrev{}𝔐}}
\newcommand\missing{{\scriptsize>}}
\newcommand\peshitta{{\abbrev{}𝔖}}
\newcommand\qumran{{\abbrev{}𝔔}}
\newcommand\sampen{{\abbrev\large ⅏}}
\newcommand\schwa{ə}
\newcommand\septuagint{{\abbrev{}𝔊}\thinspace}
\newcommand\super[1]{\textsuperscript{#1}}
\newcommand\superit[1]{\textsuperscript{\textit{#1}}}
\newcommand\symmachus{\Greek{σ´}}
\newcommand\targum{{\abbrev{}𝔗}}
\newcommand\tetragrammaton{\Hebrew{יהוה}}
\newcommand\theodotion{\Greek{θ´}}
\newcommand\vulgate{{\abbrev{}𝔙}}
% Others: http://scripts.sil.org/cms/scripts/render_download.php?format=file&media_id=ApparatusSIL_ViewGlyph&filename=ApparatusSIL_ViewGlyph.pdf

% %%%%%%%%%%%%%%%%%
%     FOOTNOTES
% %%%%%%%%%%%%%%%%%
\usepackage[perpage,multiple,hang,flushmargin]{footmisc}
% perpage: Resets footnote numbering for each page.
% multiple: Allows footnotes to be separated by a comma.
% hang and flushmargin: Removes the indentation.

% Makes footnotes lowercase, italic alpha instead of Arabic numerals.
% \renewcommand{\thefootnote}{\textit\alph{footnote}}

% Force footnotes to stay on the same page and not bleed over. Long footnotes should be placed in the appendix.
% xxxx: this isn't perfect in preventing footnotes from appearing on the next page.
\interfootnotelinepenalty=10000

% Abbreviated footnote commands
\newcommand\alt[1]{\footnote{\textsc{alt} #1}}
\newcommand\ca[2]{\footnote{\textsc{ca} #1 (#2)}}
\newcommand\cf[1]{\footnote{cf.\ #1}}
\newcommand\ed[1]{\footnote{\textsc{ed} #1}}
\newcommand\halot[2]{\footnote{\textsc{halot} \Hebrew{#1}: #2}}
\newcommand\halotu[2]{\footnote{\textsc{halotu} \Hebrew{#1}: #2}}
\newcommand\ie[1]{\footnote{\textsc{ie} #1}}
\newcommand\lit[1]{\footnote{\textsc{lit} #1}}
% \newcommand\understood{\footnote{\textsc{und}}}
\newcommand\understood[1]{\footnotemark#1\addtocounter{footnote}{-1}\footnote{\textsc{und}}}
\newcommand\ylt[1]{\footnote{\textsc{ylt} #1}}

% %%%%%%%%%%%%%%%%%%%%%%%%%
%     HEADER and FOOTER
% %%%%%%%%%%%%%%%%%%%%%%%%%
% Header options for dictionary style headers.
\usepackage{fancyhdr}

\renewcommand{\headrulewidth}{0pt} % Removes horizontal line.

\usepackage{fix2col} % Fixes potential problems with marks on two column pages.
\usepackage{substr}

% Extracts book name and chapters from rightmark and leftmark.
\fancyhead[LE]{\rightmark}
\fancyhead[RO]{\leftmark}
\fancyhead[LO,RE]{\thepage}
\fancyhead[CE]{\englishchaptertitle} % Puts book name (e.g., Genesis, Exodus) in center header on even pages.
\fancyhead[CO]{\hebrewchaptertitle} % Puts Hebrew book name in center header on odd pages.
\cfoot{} % Removes page number from center foot.

% %%%%%%%%%%%%%%
%     POETRY
% %%%%%%%%%%%%%%
% Use `\p[a-e]` for beginning of line spacing.
% Use `\pa{}` for interline separations.
\newcommand\pa{\noindent\hangindent=2em\hspace{2em}}
\newcommand\pb{\noindent\hangindent=2.5em\hspace{2.5em}}
\newcommand\pc{\noindent\hangindent=3em\hspace{3em}}
\newcommand\pd{\noindent\hangindent=3.5em\hspace{3.5em}}
\newcommand\pe{\noindent\hangindent=4em\hspace{4em}}

% Inter-line horizontal spacing
\newcommand\ilsp{\hspace{2.5em}}

% %%%%%%%%%%%%%%%%%%%%%%%%%
%     TABLE OF CONTENTS
% %%%%%%%%%%%%%%%%%%%%%%%%%
\AtBeginDocument{\addtocontents{toc}{\protect\thispagestyle{empty}}} % Removes page numbering from ToC.
\renewcommand\contentsname{CONTENTS} % Changes title of page.
\setcounter{tocdepth}{0} % Forces ToC to only displays books.
\def\numberline#1{} % Removes numbering from ToC but not from chapters.

% %%%%%%%%%%%%
%     YEAR
% %%%%%%%%%%%%
\usepackage{datetime}
\newdateformat{justtheyear}{\THEYEAR}

% %%%%%%%%%%%%%%%%%%%%%%%%%%%%%%%%%%
%     PDF BOOKMARKS AND LINKS
% %%%%%%%%%%%%%%%%%%%%%%%%%%%%%%%%%%
% Loaded last, as hyperref expects. Bookmarks list each book and its chapters;
% the table of contents and \ref links are clickable. hidelinks keeps the page
% free of colored boxes, so the PDF prints exactly as before. hyperfootnotes is
% off because per-page footnote numbering and \understood's manual \footnotemark
% would otherwise create duplicate link targets.
\usepackage[implicit=false, hidelinks, hyperfootnotes=false, bookmarksopen=false, unicode]{hyperref}
\usepackage{bookmark}
\hypersetup{
  pdftitle={The Holy Bible},
  pdfsubject={Translated from the Biblia Hebraica Stuttgartensia},
  pdfauthor={Colby Goettel},
  pdflang={en-US},
}
